\documentclass[titlepage,a4paper]{jsarticle}
\usepackage{../../../sty/import}% 各種パッケージインポート
\usepackage{../../../sty/title}% タイトルページの変更

%% タイトルページの変数
% レポートタイトル
\title{原子力材料と核燃料}
% 提出日
\expdate{\today}
% 科目名
\subject{原子力材料と核燃料}
% 分野
\class{情報経営システム工学分野}
% 学年
\grade{M1}
% 学籍番号
\mynumber{24336488}
% 記述者
\author{本間三暉}

\begin{document}
% titleページ作成
\maketitle

結晶材料に存在する欠陥は，欠陥が広がる次元によって，点欠陥，線欠陥，面欠陥，体積欠陥の4種類に分類される．

\section{点欠陥（0次元欠陥）}

点欠陥は，結晶格子中の局所的な原子配列の乱れである．代表的な点欠陥には，以下のものがある．

\begin{itemize}
    \item \textbf{空孔}\\
    本来原子が存在する格子点から原子が抜けた欠陥である．

    \item \textbf{自己格子間原子}\\
    材料自身の原子が，本来の格子点ではない格子間位置に入り込んだ欠陥である．

    \item \textbf{置換型不純物原子}\\
    母材原子の格子点を，異なる元素の原子が置き換えた欠陥である．

    \item \textbf{侵入型不純物原子}\\
    異なる元素の原子が，母材原子の格子間位置に入り込んだ欠陥である．
\end{itemize}

空孔や自己格子間原子は，材料自身の原子によって生じるため内因性点欠陥に分類される．一方，置換型不純物原子や侵入型不純物原子は，異種元素によって生じるため外因性点欠陥に分類される．

\section{線欠陥（1次元欠陥）}

線欠陥は，結晶中に線状に存在する原子配列の乱れであり，主に転位を指す．代表的な転位には，以下のものがある．

\begin{itemize}
    \item \textbf{刃状転位}\\
    結晶中に余分な原子面が入り込んだような構造を持つ転位である．

    \item \textbf{らせん転位}\\
    結晶面が転位線を中心として，らせん状にずれた構造を持つ転位である．
\end{itemize}

材料に外力が加わると，転位が結晶中を移動することで原子面のすべりが発生する．これによって，材料には元の形状に戻らない塑性変形が生じる．

\section{面欠陥（2次元欠陥）}

面欠陥は，結晶中に面状に広がって存在する欠陥である．代表例として，結晶粒界が挙げられる．

結晶粒界は，結晶方位の異なる結晶粒同士の境界である．粒界付近では，原子が結晶内部のように規則正しく配列していない．そのため，結晶粒界は材料の変形，強度，腐食などの性質に影響を与える．

\section{体積欠陥（3次元欠陥）}

体積欠陥は，三次元的な大きさを持つ欠陥である．代表的な体積欠陥には，以下のものがある．

\begin{itemize}
    \item \textbf{析出物}\\
    母相とは異なる組成や結晶構造を持つ物質が，材料内部に形成されたものである．

    \item \textbf{き裂}\\
    材料内部または表面に生じた割れであり，材料の破壊につながる場合がある．

    \item \textbf{気孔・空洞}\\
    材料内部に生じた空間であり，ポロシティとも呼ばれる．
\end{itemize}

これらの体積欠陥は，材料の強度，延性および破壊特性に大きな影響を与える．

\section{まとめ}

結晶欠陥は，欠陥の広がる次元によって，0次元の点欠陥，1次元の線欠陥，2次元の面欠陥，3次元の体積欠陥に分類される．これらの欠陥は，材料の変形，強度，延性および破壊などの機械的性質に影響を与える重要な要因である．


% 参考文献
\begin{thebibliography}{99}
\bibitem{} 原子力材料と核燃料 講義資料
\end{thebibliography}

\end{document}